\documentclass[11pt]{article}
\usepackage[margin=1in]{geometry}
\usepackage{graphicx}
\usepackage{booktabs}
\usepackage{hyperref}
\usepackage{amsmath}
\usepackage{xcolor}

\title{Exact Recovery of PDN Grid Specifications from Layout Geometry:\\
A Reverse-Engineering Study on OpenROAD's \texttt{pdngen}}
\author{Feilian Huang\\ \textit{Independent Researcher}\\ \texttt{fffeilian@gmail.com}}
\date{September 26, 2026}

\begin{document}
\maketitle

\begin{abstract}
Power delivery network (PDN) grid specifications---rail widths, stripe
pitches and offsets, layer-to-layer connectivity, and net phasing---are
normally treated as proprietary design-rule inputs to physical design.
We demonstrate that, for grids generated by OpenROAD's \texttt{pdngen},
these specifications can be recovered \emph{exactly} from the extracted
layout database.  We first characterize the complete forward model of
\texttt{pdngen} (verified against \texttt{PdnGen.tcl} at OpenROAD commit
\texttt{f12e2f47}), including its non-obvious stripe-offset reference
frame and the two distinct mechanisms that phase the ground grid relative
to the power grid.  We then implement a deterministic inference algorithm
that recovers the minimal grid specification from extracted power/ground
geometry.  On four designs across three open PDKs
(sky130hd, nangate45, asap7), \emph{every} geometry-determining parameter
(width, pitch, offset, spacing, connect pairs, \texttt{stripes\_start\_with})
is recovered exactly, down to the layout database's unit grid.
Re-running
\texttt{pdngen} with the inferred configurations reproduces the original
geometry wire-for-wire and via-for-via (4/4 designs; up to 318 wires and
16,403 vias per design).  A re-run on current OpenROAD (HEAD
\texttt{80c6be9}, September 2026) with the new C++ PDN API confirms the
mechanism is not a legacy artifact: 3,051/3,051 shapes (134 wires $+$
2,917 vias) recovered exactly on \texttt{sky130hd/gcd}.  An identifiability analysis shows the parameter-to-geometry mapping to be
empirically identifiable on the live-parameter quotient space:
\texttt{rails\_start\_with} is provably dead (never read in
\texttt{PdnGen.tcl}); rail pitch/offset are dead by code inspection plus
case verification; and four explicitly characterized degenerate
regimes---including a zero-yield \texttt{connect} pair geometrically
equivalent to omitting it, and $\mu$m-to-dbu rounding collisions---are
the only non-trivial equivalence cases observed.  We then design and evaluate anti-RE PDN
obfuscations: an interleaved decoy grid (36 extra via'd stripes at
$+P/4$) crashes the naive inference, but an adaptive two-sequence
attacker recovers all 73 stripe centers exactly---so the defense protects
\emph{intent attribution} (two exactly equivalent 9-stripe VSS grids
remain indistinguishable as to which is original), not geometry cloning;
the $-61.2\%$ IR improvement over baseline is fully explained by doubled
metal (against an equal-metal $P/2$ baseline the decoy grid is $+6.0\%$
worse in worst-case drop, at $2\times$ routing blockage and $2.06\times$
via count); pitch jitter breaks naive inference at $\geq 2\%$ but
only degrades precision (a robust least-squares attacker recovers pitch
within 1.2\% at 15\% jitter); via dropout is falsified as a defense
(\texttt{connect} inference unaffected at up to 50\% removal while IR
drop degrades by up to $+426\%$).  The result is both a practical
reverse-engineering capability and a caution: regular, rule-generated PDNs
leak their complete construction rules through the layout
database---unless the grid
is deliberately engineered not to.
\end{abstract}

\section{Introduction and Threat Model}

Hardware reverse engineering (RE) classically proceeds from silicon
upward: delayering, imaging, circuit extraction, and netlist
reconstruction~\cite{torrance2009,quadir2016,fyrbiak2017,botero2021}.  Most RE
research targets logic functionality, while the power delivery network is
treated as background infrastructure.  We show that the PDN deserves
attention in its own right: the \emph{rules} that generate it can be
recovered exactly from the extracted layout database.

\paragraph{Threat model.} The adversary obtains the physical layout
database of a placed-and-routed design (DEF/GDSII/OASIS, e.g.\ from an
untrusted foundry or a leaked tapeout database).  Concretely, our
extractor consumes the post-extraction ODB view: per-shape records carry
the power/ground net identity (from DEF \texttt{SPECIALNETS}
\texttt{USE POWER/GROUND}), the shape kind (\texttt{FOLLOWPIN} vs.\
\texttt{STRIPE}), and via bottom/top layer spans, plus die/core/row
metadata.  No access to the original PDN configuration scripts,
PDK-private design rules beyond the public LEF, or the EDA tool's
internal state is assumed.  The adversary's goal is to recover the PDN
\emph{generation} specification---the rule set that, fed back into the
same generator, reproduces the grid---rather than merely tracing wires.
An attacker starting from raw imaged polygons alone would additionally
need connectivity tracing and power/ground net identification before our
inference could run; that harder, label-free setting is not evaluated
here, and we list the unlabeled-polygon ablation as explicit future work.

The spec matters even though the geometry can always be cloned, for two
reasons.  First, \emph{regeneration under modification}: any post-silicon
tampering that must keep the power grid consistent---inserting a hardware
Trojan, re-spinning a block, or porting the layout to a derivative
design---needs the generation rules, not a frozen snapshot of shapes.
The spec is the portable, reusable form of the design IP; the geometry is
one instantiation of it.  Second, and more narrowly, \emph{partial
observation}: an attacker who recovers the rules from a sampled die
region can predict the PDN across the full chip, including areas never
imaged---something static geometry cloning cannot do.  We quantify the
limits of this in \S\ref{sec:crop}: on sky130hd/gcd, a 20\%-area strip
crop or a 10\%-area square crop recovers every parameter exactly and
regenerates the full die, while a thin 10\% strip fails for lack of met4
samples.  Recovery from a region therefore depends on the crop shape and
on each (layer, net) presenting enough period instances; and we stress
these are labeled-database crops---the harder label-free imaged-polygon
setting remains future work.

We instantiate this attack against OpenROAD's \texttt{pdngen}, the
power-grid generator of the open-source OpenROAD
flow~\cite{ajayi2019toward,ajayi2019openroad}.  We chose
\texttt{pdngen} because its source is public, which lets us verify our
forward model line-by-line, and because its rule language
(\texttt{pdngen::specify\_grid}) is representative of industrial PDN
specification styles (rails, stripes, rings, via
connectivity~\cite{kahng2011vlsi}).

\paragraph{Contributions.}
\begin{enumerate}
\item A verified forward model of \texttt{pdngen} grid generation,
including the stripe-offset reference frame and the two VSS phasing
mechanisms.
\item A deterministic inference algorithm recovering the minimal
geometry-determining grid specification from extracted PG geometry.
\item Exact recovery on 4 designs $\times$ 3 PDKs (every
geometry-determining parameter recovered exactly), with round-trip
verification:
re-generated geometry matches the original wire-for-wire and via-for-via,
plus a partial-region study quantifying when a local crop suffices for
full-die regeneration.
\item An honest identifiability analysis: the parameter-to-geometry map
is empirically identifiable on the live-parameter quotient space, with four
explicitly characterized degenerate regimes; a plausible equivalence
hypothesis is experimentally falsified.
\item An evaluation of anti-RE PDN obfuscations (interleaved decoy grids,
pitch jitter, via dropout) with attacker/defender-side measurements,
including falsified defenses and honestly reported negative controls.
\end{enumerate}

\section{Forward Model of \texttt{pdngen}}

All claims in this section were verified against \texttt{PdnGen.tcl} at
OpenROAD commit \texttt{f12e2f47} (March 2022 build).  We describe the
\texttt{stdcell} grid; macro grids follow the same stripe machinery.

\paragraph{Rails.} One rail is generated per standard-cell row boundary,
with the configured width, spanning the row width.  Net assignment
(VDD/VSS) is determined by row orientation (R0 rows: VDD on top;
MX rows: flipped---observed on our sky130hd test designs; this follows
the PDK's standard-cell definitions).  Notably, the configuration parameter
\texttt{rails\_start\_with} \emph{appears zero times} in
\texttt{PdnGen.tcl}---it is never read and is therefore a dead
parameter.  Rail \texttt{pitch}/\texttt{offset}, when present, do not
affect rail placement either (a missing pitch defaults to twice the row
height in \texttt{check\_rails}).

\paragraph{Straps.} Let the stdcell area be the union of standard-cell
rows.  \texttt{pdngen} works on an expanded area
(\texttt{stdcell\_plus\_area}): identical in $x$, but extended in $y$ by
half the maximum rail width on each side (the source comment reads
``larger than the stdcell area by half a rail'').  Stripe centers are
computed relative to the reference point
\begin{equation}
  \mathit{ref} = (\mathit{stdcell\_xMin},\;
  \mathit{stdcell\_yMin} - \mathit{max\_rail\_width}/2).
\end{equation}
The base grid (the net named by \texttt{stripes\_start\_with}, POWER
$\rightarrow$ VDD) is placed at $\mathit{ref} + \mathit{offset} + k\cdot
\mathit{pitch}$ for integer $k \geq 0$ while the center is below
$\mathit{area\_max} - \mathit{width}$.  The other net's grid is the same
grid shifted by
\begin{equation}
  \mathit{shift} =
  \begin{cases}
    \mathit{pitch}/2, & \text{no \texttt{spacing} given},\\
    \mathit{spacing} + \mathit{width}, & \text{otherwise}.
  \end{cases}
\end{equation}
Thus VDD and VSS each form a complete pitch grid; they are not a single
interleaved grid.  We verified the reference frame numerically, e.g.\
sky130hd/gcd: first VDD met4 stripe center $23690 = 10120 + 13570$ dbu,
and first VDD met5 stripe center $24240 = (10880 - 240) + 13600$ dbu,
where $240$ is half the $480$\,dbu rail width.

\paragraph{Vias and \texttt{connect}.} \texttt{connect} declares logical
layer pairs, e.g.\ \texttt{\{\{met1 met4\} \{met4 met5\}\}}.
\texttt{pdngen} realizes each pair as a via stack spanning all
intermediate layers, placed at same-net wire crossings: a pair may span
several routing layers---asap7's \texttt{\{M1 M5\}} generates a
continuous four-segment stack (M1$\to$M2$\to$M3$\to$M4$\to$M5, 106 vias
per segment).  Whether insertion succeeds is decided not by layer
``adjacency'' but by the geometric overlap of the two endpoint grids at
each crossing: orthogonal (H$\times$V) crossings succeed, while a
parallel narrow rail under a wide stripe, whose overlap cannot cover the
stripe's full width, is skipped entirely with PDN-0042 warnings---which
is why sky130hd's \texttt{\{met1 met5\}} (both horizontal) yields
\emph{zero} vias.  We return to this in Section~\ref{sec:ident}.

\section{Inference Algorithm}

Input: extracted PG geometry (per-shape records: net, layer, shape,
kind $\in$ \{wire, via\}, bounding box; via records carry the bottom/top
layer span), plus die/core/row metadata and the DEF \texttt{SPECIALNETS}
section (to read power/ground net names from \texttt{USE POWER/GROUND}).
The ground-truth configuration is never consulted.  Note this is the
labeled layout-database view described in the threat model---net identity
and shape kinds come from the ODB/DEF metadata, not from unlabeled
polygons.

\paragraph{Rails.} \texttt{FOLLOWPIN} wires are grouped by layer; width is
the modal cross-dimension.  Rail positions are row-determined and need no
further parameters.

\paragraph{Straps.} \texttt{STRIPE} wires are grouped by (net, layer).
Per layer and net, sorted stripe centers are fit by least squares to
$c_i = c_0 + i\cdot p$ (pitch $p$ rounded to the unit grid; residual gate
2\,dbu, which rejects genuinely multimodal inputs loudly rather than
silently accepting them).  The
net with the smaller first center is the base grid, fixing
\texttt{stripes\_start\_with} (POWER or GROUND).  Offset is
$(\text{first base center} - \mathit{ref})$ with $\mathit{ref}$ from
Eq.~(1).  The VSS shift is compared against $\mathit{pitch}/2$: equality
(no \texttt{spacing} emitted) versus $\mathit{spacing} = \mathit{shift} -
\mathit{width}$.

\paragraph{\texttt{connect}.} Via (bottom, top) spans are chained
bottom-to-top through shared layers.  Within each chain, the lowest
\emph{endpoint} rail layer paired with the first strap endpoint above it
yields the rail$\to$strap pair; consecutive strap endpoints within the
chain yield strap$\to$strap pairs.  Intermediate layers of a stack are not
endpoints: in asap7 the chain M1$\to$M2$\to$M3$\to$M4$\to$M5$\to$M6
collapses to \texttt{\{\{M1 M5\} \{M5 M6\}\}}---M2, though a rail layer,
is a middle layer of the M1$\to$M5 stack, so the inferred pair is
M1$\to$M5, matching ground truth.

The output is a minimal \texttt{pdngen::specify\_grid} configuration
containing only geometry-determining parameters; dead parameters
(\texttt{halo}, \texttt{rails\_start\_with}, rail pitch/offset, grid name)
are omitted.

\section{Experimental Setup}

We built OpenROAD (\texttt{f12e2f47}, March 2022) and Yosys 0.69
(oss-cad-suite) via conda, with OpenROAD-flow-scripts at commit
\texttt{96eb3de} (April 2022).  Yosys is used only for synthesis; the PDN
geometry under study is independent of the Yosys version.  Four designs were run through floorplan and PDN generation:
\texttt{sky130hd/gcd}, \texttt{nangate45/gcd}, \texttt{sky130hd/aes},
and \texttt{asap7/gcd}.  This ORFS vintage uses the legacy
\texttt{pdngen::specify\_grid} syntax (via \texttt{pdn.cfg} /
\texttt{grid\_strategy} files), which is our inference target.
Geometry was extracted directly from the layout database with the Tcl ODB
API (\texttt{getSWires}/\texttt{getWires} over POWER/GROUND nets),
because this build's \texttt{write\_def} degrades stripe wires to
zero-width points.  None of the four designs contains a core ring, so the
evaluation covers rails, stripes, and via stacks.

\section{Results}

\subsection{Exact parameter recovery}

Table~\ref{tab:recovery} compares inferred versus ground-truth
parameters.  \emph{Every} geometry-determining parameter matches exactly,
down to the layout database's unit grid (width, pitch, offset, spacing;
connect pairs and \texttt{stripes\_start\_with} identically).  The asap7 case is notable:
its ground truth selects between two grid strategies by a conditional, and
inference correctly identified the active branch, including the
\texttt{spacing}-based (rather than pitch/2) VSS phasing on both strap
layers (M5 shift $192\,\mathrm{dbu} = 72 + 120$; M6 shift
$384\,\mathrm{dbu} = 96 + 288$).

\begin{table}[h]
\centering
\small
\caption{Inferred vs.\ ground-truth grid parameters (all match exactly).}
\label{tab:recovery}
\begin{tabular}{lp{2.0cm}p{6.0cm}p{4.4cm}}
\toprule
Design & Rails & Straps (width / pitch / offset [/ spacing]) & \texttt{connect} \\
\midrule
sky130hd/gcd &
met1 0.480 &
met4 1.600/27.140/13.570; met5 1.600/27.200/13.600 &
\{\{met1 met4\} \{met4 met5\}\} \\
nangate45/gcd &
metal1 0.170 &
metal4 0.480/56.000/2.000; metal7 1.400/40.000/2.000 &
\{\{metal1 metal4\} \{metal4 metal7\}\} \\
sky130hd/aes &
met1 0.480 &
(same as sky130hd/gcd) &
\{\{met1 met4\} \{met4 met5\}\} \\
asap7/gcd &
M1, M2 0.018 &
M5 0.120/11.880/0.300/sp.\,0.072; M6 0.288/12.000/0.513/sp.\,0.096 &
\{\{M1 M5\} \{M5 M6\}\} \\
\bottomrule
\end{tabular}
\end{table}

\subsection{Round-trip verification: 4/4 exact}

We re-ran the ORFS PDN step with only the inferred configuration
(isolated output directory; inputs symlinked), re-extracted geometry, and
compared wire rectangles (1\,dbu tolerance) and via sets.  All four
designs reproduce the original geometry \emph{exactly}:

\begin{center}
\begin{tabular}{lrrc}
\toprule
Design & Wires (orig / regen) & Vias (orig / regen) & Result \\
\midrule
sky130hd/gcd & 133 / 133 & 2,907 / 2,907 & PASS \\
nangate45/gcd & 65 / 65 & 279 / 279 & PASS \\
sky130hd/aes & 318 / 318 & 16,403 / 16,403 & PASS \\
asap7/gcd & 114 / 114 & 432 / 432 & PASS \\
\bottomrule
\end{tabular}
\end{center}

Because the inferred configurations omit \texttt{halo},
\texttt{rails\_start\_with}, and rail pitch/offset, the round trip also
\emph{empirically} proves these parameters have no geometric effect.

\subsection{Variant experiments (sensitivity and dead parameters)}

On sky130hd/gcd we additionally verified: (i) changing met4 pitch
$27.14 \to 30.00$ changes geometry (2,907 $\to$ 2,601 vias) and
re-inference recovers $30.000$ exactly; (ii) flipping
\texttt{stripes\_start\_with} to GROUND swaps the VDD/VSS grids and is
recovered as GROUND; (iii) adding \texttt{rails\_start\_with}=GROUND,
\texttt{halo}=5, and rail pitch/offset leaves geometry bit-identical
(133 wires / 2,907 vias)---direct evidence of dead parameters.

\subsection{Partial-region extrapolation}
\label{sec:crop}

To test the partial-observation motivation honestly, we cropped the
sky130hd/gcd extracted geometry---the labeled ODB view, not imaged
polygons---to three windows and re-ran inference on each crop alone.  A
20\%-width full-height vertical strip ($\sim$20\% die area) recovers every
parameter exactly---met4 1.600/27.140/13.570, met5 1.600/27.200/13.600,
\texttt{connect} and \texttt{stripes\_start\_with} all exact---and
regenerates
the full die: PASS.  A $\sim$10\%-area square in the lower-left corner
(3 stripes per layer per net in-window) likewise recovers everything
exactly: PASS.  But a thin 10\%-width vertical strip fails: met5 is
recovered exactly, yet the window holds only one VDD and zero VSS met4
stripes, so met4 pitch is unidentifiable and full-die regeneration
FAILS---an information deficit, not a numerical error.  Local recovery
therefore depends on crop shape and on each (layer, net) presenting
enough period instances (here, $\geq$2--3 stripes sufficed).  We stress
the gap this does not close: these are labeled-database crops; net
identification and shape-kind labeling from unlabeled imaged polygons
remain unevaluated future work.

\section{Identifiability Analysis}
\label{sec:ident}

We began this study expecting the parameter$\to$geometry mapping to be
significantly non-injective, hoping to characterize large equivalence
classes.  The experiments say otherwise.  Formally, let $G$ be the map
from \texttt{pdngen} configurations to placed grid geometry, and let
$\sim$ identify configurations that differ only in dead parameters:
\texttt{rails\_start\_with} (provably dead---it appears zero times in
\texttt{PdnGen.tcl}), \texttt{halo} (macro-only), rail
\texttt{pitch}/\texttt{offset} (dead by code inspection and case
verification: re-adding them leaves geometry bit-identical), and the
arbitrary grid name.
\textbf{On the quotient space of live parameters, $G$ is empirically
identifiable}---every live parameter was recovered exactly on all tested
configurations---\textbf{except in the following degenerate regimes},
which we state explicitly as the non-trivial equivalence cases observed:

\begin{enumerate}
\item \emph{Phasing coincidence.} When an explicit \texttt{spacing} $S$
and width $W$ satisfy $S + W = \mathit{pitch}/2$, the resulting VSS phase
coincides with the no-\texttt{spacing} default, so the explicit-spacing
and no-spacing configurations are geometrically identical.
\item \emph{Single-stripe layers.} A layer carrying exactly one stripe
leaves \texttt{pitch} unidentifiable: any pitch above a geometry-implied
lower bound reproduces the same shapes.
\item \emph{Zero-yield connect pairs.} A \texttt{connect} pair that
produces no vias is geometrically equivalent to omitting that pair
entirely.  The mechanism is geometric overlap, not layer adjacency
(\S2): the spanning pair \texttt{\{\{met1 met5\}\}} yields \emph{zero}
vias because the parallel narrow met1 rail cannot cover the wide met5
stripe---\texttt{pdngen} skips every insertion (PDN-0042)---while the
split pair \texttt{\{\{met1 met4\} \{met4 met5\}\}} yields the full
2,907-via stack, so spanning-vs-split is identifiable (connect-pair
granularity is recoverable), but spanning-vs-omitted is a genuine
equivalence.  Note this zero-yield behavior was observed on the 2022
legacy Tcl implementation; the current C++ version turns it into a hard
error (PDN-0179, see \S\ref{sec:deflim}).
\item \emph{Micron-to-dbu rounding.} Stripe geometry is computed in
database units, so distinct $\mu$m values that round to the same dbu are
indistinguishable.  Recovery is therefore exact only down to the layout
database's unit grid, and parameter values must be quoted with that
quantization in mind.
\end{enumerate}

Outside these regimes, every geometry-determining parameter was recovered
exactly from the extracted layout database on every tested
configuration.  Macro-grid parameters, pin names, and blockages are
unrecoverable only when absent from the design (no macros, hence no
macro-grid geometry to observe).  We deliberately weaken ``injective''
to \emph{empirically identifiable}: the evidence is four designs plus
variant experiments, not a proof, and joint transforms we did not test
(e.g.\ offset $\geq$ pitch/2 interacting with
\texttt{stripes\_start\_with}) could in principle add equivalence cases.
The claim is therefore a precise empirical statement about the
live-parameter quotient space, not a dismissal of the degenerate cases
above.

\section{Defenses: Anti-RE PDN Obfuscation}
\label{sec:defense}

The attack succeeds because rule-generated PDNs are perfectly regular:
uniform pitches, deterministic VSS phasing, and complete via stacks at
every crossing make the forward model invertible.  A natural counter is
\emph{anti-RE PDN obfuscation}: grids engineered so that the inference
assumptions fail, while the PDN remains functional.  Hardware obfuscation
has a rich literature at the logic level (camouflaging, logic
locking~\cite{shamsi2019,rajendran2012}), in split
manufacturing~\cite{perez2020,imeson2013}, and in design-for-trust
generally~\cite{rajendran2014}; physical-design
watermarking~\cite{kahng2001} and power-signature
watermarking~\cite{ziener2008} hide information in layout and power
networks respectively, and layout-level anti-RE
techniques~\cite{chen2015} are the closest defense-side analog to our
decoy grids.  On the attack side, our work is most directly preceded by
ReGDS~\cite{rajarathnam2020}, which recovers gate-level netlists from
GDSII at 100\%---we invert the same ``layout as evidence'' threat model
one level up, recovering generator \emph{parameters} rather than
netlists---and the power grid's security relevance is independently
established by power-grid-aware side-channel
analysis~\cite{wang2013} and PDN-impedance tamper
sensing~\cite{mosavirik2023}.  A systematic search found no peer-reviewed
prior work on obfuscation of the on-chip PDN structure itself.
(Prior metal-mesh work targets invasive probing and tamper
evidence---conductive shield layers that detect decapsulation---a
different threat model from rule extraction, which we do not count as
PDN-grid obfuscation prior art.)

\paragraph{Defense goal and methodology.} The defender's goal is to make
the grid \emph{specification}---not the geometry, which an attacker can
always clone---unrecoverable or ambiguous, while keeping the PDN
functional: fully connected, DRC-clean, and with no significant IR-drop
degradation.  Naming: we label experiments E1--E3 and the underlying
defense mechanisms D1--D4 (design-doc numbering)---D1: decoys without
vias or net attachment, D2: pitch jitter, D3: interleaved decoy grids,
D4: via dropout---so E1 evaluates D3, E2 evaluates D2, and E3 evaluates
D4.  We evaluate candidate obfuscations on sky130hd/gcd against
two attacker models: the naive deterministic inference of
Section~3 (zero modifications), and adaptive attackers that know a
defense may be present.  Defender-side checks are: (i) connectivity of
the wire/via graph per net (a single component, no dangling vias); (ii)
LEF minimum spacing; (iii) a lumped resistive mesh (LEF
\texttt{RPERSQ}, uniform 1\,A rail-current injection, ideal met5 feed,
5\,$\Omega$ nominal via resistance) used strictly as a
\emph{comparative} IR metric, never as signoff.  Experiments manipulate
the extracted geometry directly (the complete input of
\texttt{pg\_infer.py}), which is equivalent to ODB-level edits for every
attacker-side measurement; the honest gap is that signal-wire DRC around
decoys is not modeled, though decoys sit in the $+P/4$ gaps
($5.19\,\mu\mathrm{m}$) far from any DRC risk region, with LEF spacing
rules verified.

\subsection{E1: Interleaved decoy grids (effective)}

\paragraph{Setup.} On sky130hd/gcd (met4 pitch $27140$\,dbu, met5 pitch
$27200$\,dbu), every true stripe is duplicated at $+P/4$ with the same
width, the same net, and a full copy of its via stack (met1--met4 stack
plus the met4--met5 crossing vias), faithfully simulating a second
\texttt{pdngen} pass interleaved with the first; decoy$\times$decoy
same-net met4--met5 vias are added as well.  Stripes obeying
\texttt{pdngen}'s loop bound yield 36 extra decoy stripes (18 per layer):
wires $133 \to 169$, vias $2{,}907 \to 5{,}994$.

\paragraph{Attacker side.} The naive inference fails loudly: with the
least-squares fit of \S3, met4 VDD reports pitch $13570.0$\,dbu with
maximum residual $3571.1$\,dbu, far above the 2\,dbu gate
(\texttt{AssertionError}; no parameters emitted).  The exact-recovery
attack of Sections~3--5 is defeated outright---for the naive attacker.

\paragraph{Adaptive attacker: multi-sequence inference.} An attacker who
knows decoys may be present upgrades the pitch fit to $k$ arithmetic
sequences per (layer, net): $P$ from the modal adjacent difference,
phases from circular clustering of $(c \bmod P)$, each cluster validated
as an arithmetic progression.  With $k=2$ this attacker recovers
\emph{all 73 stripe centers exactly}---met4 $P = 27.140\,\mu\mathrm{m}$
with phases at truth and truth$+P/4$ (19/19 VDD, 18/18 VSS centers), met5
$P = 27.200\,\mu\mathrm{m}$ likewise (18/18 per net).  The prediction is
confirmed: both attack goals of Section~1 (extrapolation, regeneration
under modification) need only this parameter bundle, not knowledge of
which sequence is original.  D3 therefore does \emph{not} prevent
geometry cloning.\footnote{Strictly, this parameter bundle is not a
single legacy-\texttt{pdngen} configuration: feeding two
\texttt{specify\_grid stdcell} blocks silently ignores the second (one
\texttt{PDN-0013} in the log; first-writer-wins), and running
\texttt{pdngen} twice rips up the first grid.  E1a's interleaved geometry
was constructed geometrically, which is all the attacker needs---the
bundle regenerates shapes directly.}

\paragraph{What D3 does protect: intent attribution.} Even with the full
two-sequence bundle recovered, the attacker cannot determine which
sequence is the designer's original grid.  Enumerating stripe subsets on
met4 VSS (9 true + 9 decoy stripes) yields two \emph{exactly equivalent}
9-stripe pitch grids---the true grid and the decoy grid.  On VDD (10 true
+ 9 decoy), three 9-stripe candidates exist (two true 9-subsequences plus
the decoy grid); the full 10-stripe grid is unique, but identifying it
requires the prior knowledge that the original had 10 stripes, which the
attacker does not have, and the defender can invert the majority so that
``largest AP is the original'' heuristics fail.  This is \textbf{genuine
intent ambiguity}: two stripe subsets both form valid
\texttt{pdngen}-style pitch grids, and the attacker cannot determine
which was the designer's original.  (This is a different object from the
configuration-level equivalence classes hunted in
Section~\ref{sec:ident}---there, two generator configurations mapping to
the same geometry; here, two subsets of one defended geometry each
admitting a valid generator configuration.)  The honest boundary: D3
blinds the naive attacker and reduces the adaptive attacker from exact
\emph{attribution} to ambiguity; it does not impede geometry recovery at
all.

\paragraph{Defender side.} VDD and VSS each remain a single connected
component with zero dangling vias.  Minimum stripe spacing is
$5.185\,\mu\mathrm{m}$ on met4 ($\geq 0.3\,\mu\mathrm{m}$ LEF) and
$5.2\,\mu\mathrm{m}$ on met5 ($\geq 1.6\,\mu\mathrm{m}$ LEF).  The mesh
reports worst-case drop $66.90 \to 25.98$\,mV/A and mean $38.89 \to
11.63$\,mV/A---but this $-61.2\%$ is \emph{fully explained by doubled
metal}, not by obfuscation.  Against an equal-metal unobfuscated baseline
(met4/met5 pitch halved to $P/2$, no decoys; 169 wires, 5,994 vias,
stripe length $18{,}918.4\,\mu\mathrm{m}$---identical to E1a):

\begin{center}
\begin{tabular}{lrrr}
\toprule
 & Baseline & E1a (D3 decoys) & $P/2$ equal-metal \\
\midrule
wires & 133 & 169 & 169 \\
vias & 2,907 & 5,994 & 5,994 \\
stripe length ($\mu$m) & 9,588.6 & 18,918.4 & 18,918.4 \\
worst IR (mV/A) & 66.90 & 25.98 & 24.51 \\
mean IR (mV/A) & 38.89 & 11.63 & 11.54 \\
\bottomrule
\end{tabular}
\end{center}

The decoy grid is $+6.0\%$ \emph{worse} than the equal-metal baseline in
worst-case drop (25.98 vs.\ 24.51\,mV/A) and $+0.8\%$ in mean (11.63
vs.\ 11.54\,mV/A)---essentially tied.  We therefore withdraw the
``negative physical cost'' framing: D3's real price is $2\times$ routing
blockage on met4/met5 and $2.06\times$ via count, for no IR benefit over
simply halving the pitch.

\subsection{E2: Pitch jitter (degrades precision, does not blind)}

\paragraph{Setup.} Each met4 stripe and its full via stack are translated
by $\delta \sim \mathrm{Uniform}(-j, +j)$ for $j \in \{0, 2\%, 5\%, 10\%,
15\%\}$ of pitch, three seeds each.  Stripes violating LEF spacing or via
landing (vias must stay within the rail/met5 perpendicular overlap
extent) are resampled, so all 13 runs are physically legal.

\paragraph{Attacker side.} The naive inference breaks at $j \geq 2\%$
(its $\pm 2$\,dbu arithmetic-sequence assertion is far exceeded).  But a
$\sim$30-line robust least-squares pitch fit---the obvious attacker
upgrade---still recovers pitch with mean relative error $0.14\%$,
$0.42\%$, $0.40\%$, $0.50\%$ (max $0.21\%$, $0.76\%$, $0.97\%$,
$1.19\%$) at $j = 2\%, 5\%, 10\%, 15\%$, and offset within
$2.97\,\mu\mathrm{m}$ at $15\%$.  Jitter's honest ceiling is therefore
``degrades precision,'' not ``blinds'': it converts exact recovery into
approximate recovery, which may suffice for a foundry cloning the grid
but not for hiding the design intent.

\paragraph{Defender side.} All 13 runs keep a single connected component
per net, zero dangling vias, minimum spacing $\geq 5.68\,\mu\mathrm{m}$,
and mesh worst-case drop within $\pm 0.62\%$ of baseline---jitter is
electrically free, but buys little security against an adaptive attacker.

\subsection{E3: Via dropout (falsified as a defense)}

Deleting $p \in \{10\%, 30\%, 50\%\}$ of the 912 rail$\to$strap via
stacks (2,736 vias) leaves \texttt{connect} inference \emph{unchanged}
(\texttt{\{\{met1 met4\} \{met4 met5\}\}}): recovery operates on via
\emph{span types}, not instances, so instance deletion cannot move the
attacker's output---while mesh worst-case drop degrades by up to
$+426\%$.  Deleting the entire met2--met3 span does corrupt
\texttt{connect} (mis-inferred as \texttt{\{\{met4 met5\}\}}), but it
fragments the PDN into 49 components per net and disconnects all 192
rail injection nodes; an analytic projection shows the mis-inferred
configuration would regenerate with zero rail$\to$strap stacks.  Via
dropout is self-defeating: a falsified defense, kept only because it
quantifies that connect identifiability depends on via-chain integrity
(Section~\ref{sec:ident}).

\subsection{Negative controls}
\label{sec:neg}

We report three honestly weak mechanisms rather than hiding them.
\textbf{Floating decoys} (stripes with no net) never enter inference at
all: our extractor iterates only POWER/GROUND nets, so they are filtered
before inference begins.  \textbf{Via-less decoys} (same net, no vias)
crash the naive inference, but a one-line adaptive filter---dropping
stripes with no vias---restores \emph{exact} recovery of every
parameter, so this defense is falsified.  \textbf{Half-pitch ($+P/2$)
decoys} are physically unrealizable under the no-spacing phasing used
here: the VSS grid sits exactly at
$+P/2$ from VDD, so such decoys coincide with the opposite net's stripe
centers---a hard short---and were never run.

\subsection{Summary and limitations}
\label{sec:deflim}

\begin{center}
\small
\begin{tabular}{lp{2.4cm}p{4.2cm}p{4.2cm}}
\toprule
Mechanism & Naive attacker & Adaptive attacker & Physical cost \\
\midrule
D3 interleaved decoys (E1) & crashes (resid.\ 3571\,dbu) & geometry exact via 2-seq.\ fit; 2-way \emph{intent} ambiguity (VSS) & IR $+6.0\%$ worst vs.\ equal-metal; $2\times$ blockage, $2.06\times$ vias \\
D2 jitter (E2) & breaks at $\geq 2\%$ & pitch err.\ $\leq 1.2\%$ & IR $\pm 0.6\%$ \\
D4 via dropout (E3) & connect unchanged & wire params unaffected & IR $+42\%$--$+426\%$; self-defeating \\
D1 no-via decoys & crashes & exact recovery (falsified) & none \\
D1 floating decoys & no effect & --- & wasted routing \\
\bottomrule
\end{tabular}
\end{center}

Limitations must be stated plainly.  The resistive mesh is a comparative
first-order model (uniform current, ideal met5 feed, nominal
$5\,\Omega$ vias, no packaging or EM analysis)---it cannot be presented
as signoff, and no commercial IR-drop/EM signoff is available in this
environment.  The routability cost of D3 is now quantified against an equal-metal
baseline: $2\times$ routing blockage on met4/met5 and $2.06\times$ via
count, with worst-case IR $+6.0\%$ versus the unobfuscated $P/2$ grid; a
full-flow routing congestion/timing measurement remains future work.  The
main attack experiments (\S3--\S5) use a toy design (gcd) and OpenROAD
\texttt{f12e2f47} (2022) with the legacy Tcl \texttt{pdngen}
(\texttt{pdngen::specify\_grid}).  Only three distinct grid
configurations are exercised (sky130hd/gcd and aes share one), and none
of the test designs has a core ring or macro grids---the most complex,
IP-valuable part of a PDN spec.  Porting the method to commercial
generators (e.g.\ Innovus \texttt{addStripe}) or ring/macro-bearing
designs is future work.  To check that the exact-recovery
mechanism is not a legacy-only artifact, we rebuilt current OpenROAD from
source (HEAD \texttt{80c6be9}, September 2026) and re-ran the full
round-trip on \texttt{sky130hd/gcd} under the new C++ API
(\texttt{define\_pdn\_grid} / \texttt{add\_pdn\_stripe} /
\texttt{add\_pdn\_connect}), using the 2022 ORFS sky130hd LEF as input
data.  The inferred configuration regenerates the truth geometry
\emph{exactly}: 3,051/3,051 shapes (134 wires $+$ 2,917 vias), zero
symmetric difference, with met4/met5 widths, pitches, and offsets
recovered to the unit grid (met4: 1.600/27.140/13.570; met5:
1.600/27.200/13.600).  The small count difference versus the 2022 run
(133 wires / 2,907 vias) comes from the changed stripe-offset datum: the
legacy reference subtracts half the maximum rail width (Eq.~1), while the
current version references the core-area edge directly, shifting one
stripe in or out of the loop bound.  A successful \texttt{\{met1 met4\}} connect now
yields a \emph{complete} intermediate via stack (912 vias per
met1--met2, met2--met3, and met3--met4 step, 181 for met4--met5).
Generator quirks differ from the 2022 version in detail: the
non-adjacent \texttt{\{met1 met5\}} connect is now a hard error
([ERROR PDN-0179] ``unable to repair all channels'') rather than a
silent zero-via skip; \texttt{-verbose} and \texttt{-pin\_direction}
are explicitly marked deprecated (parsed but with no geometric effect);
\texttt{add\_pdn\_stripe -width 0} is still rejected (PDN-0106, minimum
width $0.3000\,\mu\mathrm{m}$).  The current-version re-run therefore
confirms the forward model and the exact-recovery mechanism, while the
defense experiments (E1--E3) remain validated on the 2022 version only.
Finally, these defenses protect the grid \emph{specification} (design IP)
against rule extraction; they do not stop an attacker who simply clones
the geometry---a different threat model.  We also note the attack surface
widens on commercial generators and on designs with core rings and macro
grids, neither of which our test set covers.

\section{Conclusion}

We demonstrated exact, automated recovery of PDN grid generation rules
from the extracted layout database, validated end-to-end by round-trip
regeneration on four designs across three PDKs, and confirmed not to be
a legacy-only artifact by an exact round-trip on current OpenROAD
(HEAD \texttt{80c6be9}) with the C++ PDN API.  The parameter-to-geometry
mapping of \texttt{pdngen} proved empirically identifiable on the
live-parameter quotient space, with four explicitly characterized
degenerate regimes---an honest, somewhat surprising negative result for
the equivalence-class hunt.  For defenders, the takeaway is
that regular PDN grids leak their complete construction rules; our defense
experiments show this leakage can be engineered into \emph{intent}
ambiguity---interleaved decoy grids leave an adaptive attacker unable to
tell which of two equivalent grids is original---but not into geometry
unclonability, and not for free (no IR benefit over an equal-metal
unobfuscated grid); cheap fixes (jitter, via dropout) provably
do not suffice.  For the open-source EDA community, the verified forward
model and the dead-parameter findings (notably
\texttt{rails\_start\_with}) are directly actionable.

\bibliographystyle{plain}
\begin{thebibliography}{18}

\bibitem{ajayi2019toward}
T.~Ajayi, V.~A. Chhabria, M.~Foga{\c{c}}a, S.~Hashemi, A.~Hosny,
A.~B. Kahng, M.~Kim, J.~Lee, U.~Mallappa, M.~Neseem, et~al.
\newblock Toward an open-source digital flow: First learnings from the
  openroad project.
\newblock In {\em Proceedings of the 56th Annual Design Automation
  Conference (DAC)}, pages 1--4, 2019.

\bibitem{ajayi2019openroad}
T.~Ajayi, D.~Blaauw, T.~B. Chan, C.~K. Cheng, V.~A. Chhabria, D.~K. Choo,
M.~Coltella, S.~Dobre, R.~Dreslinski, M.~Foga{\c{c}}a, et~al.
\newblock OpenROAD: Toward a self-driving, open-source digital layout
  implementation tool chain.
\newblock In {\em Proc.\ GOMACTECH}, pages 1105--1110, 2019.

\bibitem{kahng2011vlsi}
A.~B. Kahng, J.~Lienig, I.~L. Markov, and J.~Hu.
\newblock {\em VLSI Physical Design: From Graph Partitioning to Timing
  Closure}.
\newblock Springer, 2011.

\bibitem{torrance2009}
R.~Torrance and D.~James.
\newblock The state-of-the-art in IC reverse engineering.
\newblock In {\em Cryptographic Hardware and Embedded Systems---CHES 2009},
  volume 5747 of {\em Lecture Notes in Computer Science}, pages 363--381.
  Springer, 2009.

\bibitem{fyrbiak2017}
M.~Fyrbiak, S.~Strau{\ss}, C.~Kison, S.~Wallat, M.~Elson, N.~Rummel, and
  C.~Paar.
\newblock Hardware reverse engineering: Overview and open challenges.
\newblock In {\em 2017 IEEE 2nd International Verification and Security
  Workshop (IVSW)}, pages 88--94. IEEE, 2017.

\bibitem{quadir2016}
S.~E. Quadir, J.~Chen, D.~Forte, N.~Asadizanjani, S.~Shahbazmohamadi,
  L.~Wang, J.~Chandy, and M.~Tehranipoor.
\newblock A survey on chip to system reverse engineering.
\newblock {\em ACM Journal on Emerging Technologies in Computing Systems},
  13(1), Article~6, pages 6:1--6:34, 2016.

\bibitem{perez2020}
T.~D. Perez and S.~Pagliarini.
\newblock A survey on split manufacturing: Attacks, defenses, and
  challenges.
\newblock {\em IEEE Access}, volume 8, pages 184013--184035, 2020.

\bibitem{shamsi2019}
K.~Shamsi, M.~Li, K.~Plaks, S.~Fazzari, D.~Z. Pan, and Y.~Jin.
\newblock IP protection and supply chain security through logic
  obfuscation: A systematic overview.
\newblock {\em ACM Transactions on Design Automation of Electronic
  Systems}, 24(6):65:1--65:36, 2019.

\bibitem{botero2021}
U.~J. Botero, R.~Wilson, H.~Lu, M.~T. Rahman, M.~A. Mallaiyan, F.~Ganji,
  N.~Asadizanjani, M.~M. Tehranipoor, D.~Woodard, and D.~Forte.
\newblock Hardware trust and assurance through reverse engineering: A
  tutorial and outlook from image analysis and machine learning
  perspectives.
\newblock {\em ACM Journal on Emerging Technologies in Computing Systems},
  17(4):1--53, 2021.

\bibitem{rajarathnam2020}
R.~S. Rajarathnam, Y.~Lin, Y.~Jin, and D.~Z. Pan.
\newblock ReGDS: A reverse engineering framework from GDSII to gate-level
  netlist.
\newblock In {\em Proc.\ IEEE International Symposium on Hardware Oriented
  Security and Trust (HOST)}, pages 154--163, 2020.

\bibitem{chen2015}
S.~Chen, J.~Chen, D.~Forte, J.~Di, M.~Tehranipoor, and L.~Wang.
\newblock Chip-level anti-reverse engineering using transformable
  interconnects.
\newblock In {\em Proc.\ IEEE International Symposium on Defect and Fault
  Tolerance in VLSI and Nanotechnology Systems (DFTS)}, pages 109--114,
  2015.

\bibitem{wang2013}
X.~Wang, W.~Yueh, D.~B. Roy, S.~Narasimhan, Y.~Zheng, S.~Mukhopadhyay,
  D.~Mukhopadhyay, and S.~Bhunia.
\newblock Role of power grid in side channel attack and power-grid-aware
  secure design.
\newblock In {\em Proc.\ 50th ACM/EDAC/IEEE Design Automation Conference
  (DAC)}, 2013.

\bibitem{mosavirik2023}
T.~Mosavirik, S.~K. Monfared, M.~Saadat-Safa, and S.~Tajik.
\newblock Silicon echoes: Non-invasive Trojan and tamper detection using
  frequency-selective impedance analysis.
\newblock {\em IACR Transactions on Cryptographic Hardware and Embedded
  Systems}, 2023(4):238--261, 2023.

\bibitem{rajendran2014}
J.~Rajendran, V.~Jyothi, O.~Sinanoglu, and R.~Karri.
\newblock Regaining trust in VLSI design: Design-for-trust techniques.
\newblock {\em Proceedings of the IEEE}, 102(8):1266--1282, 2014.

\bibitem{imeson2013}
F.~Imeson, A.~Emtenan, S.~Garg, and M.~Tripunitara.
\newblock Securing computer hardware using 3D integrated circuit (IC)
  technology and split manufacturing for obfuscation.
\newblock In {\em Proc.\ 22nd USENIX Security Symposium}, pages 495--510,
  2013.

\bibitem{rajendran2012}
J.~Rajendran, Y.~Pino, O.~Sinanoglu, and R.~Karri.
\newblock Security analysis of logic obfuscation.
\newblock In {\em Proc.\ 49th ACM/EDAC/IEEE Design Automation Conference
  (DAC)}, pages 83--89, 2012.

\bibitem{kahng2001}
A.~B. Kahng, J.~Lach, W.~H. Mangione-Smith, S.~Mantik, I.~L. Markov,
  M.~Potkonjak, P.~Tucker, H.~Wang, and G.~Wolfe.
\newblock Constraint-based watermarking techniques for design IP
  protection.
\newblock {\em IEEE Transactions on Computer-Aided Design of Integrated
  Circuits and Systems}, 20(10):1236--1252, 2001.

\bibitem{ziener2008}
D.~Ziener and J.~Teich.
\newblock Power signature watermarking of IP cores for FPGAs.
\newblock {\em Journal of Signal Processing Systems}, 51(1):123--136,
  2008.

\end{thebibliography}

\end{document}
